% ----------------------------------------------------------
% Revisão da Literatura
% 
% ----------------------------------------------------------


\chapter[Referencial Teórico]{Referencial Teórico}

{\section{\textbf{A Metabolômica}}}

A metabolômica consiste no estudo sistemático das ``impressões digitais'' químicas únicas deixadas por processos celulares, especificamente os perfis de metabólitos de moléculas pequenas (geralmente <1500 Da) presentes em células, tecidos e fluidos biológicos \cite{LILAND2011827,picone2026integrating}. Considerada o nível final da hierarquia das ciências ômicas, vindo após a genômica, transcriptômica e proteômica, ela fornece um instantâneo direto do estado fisiológico e patológico em tempo real, servindo como o elo mais próximo entre o genótipo e o fenótipo \cite{picone2026integrating,LILAND2011827,anwardeen2023statistical}. Essa abordagem é amplamente utilizada em aplicações como descoberta de biomarcadores, toxicologia, nutrição, medicina de precisão, agricultura e biologia de sistemas.

Com o avanço de tecnologias analíticas, tornou-se possível quantificar simultaneamente centenas ou milhares de metabólitos \cite{li2016performance}. As duas plataformas principais são a Espectrometria de Massas (MS) e a Ressonância Magnética Nuclear (NMR) \cite{anwardeen2023statistical}. A espectroscopia de NMR é uma técnica não destrutiva, altamente reprodutível e que exige preparação mínima de amostras, embora possua sensibilidade limitada que pode mascarar compostos em baixas concentrações \cite{mayonu2026evaluation}. Já a MS, frequentemente acoplada a técnicas de separação como cromatografia líquida (LC-MS) ou gasosa (GC-MS), oferece elevada sensibilidade e seletividade, permitindo a detecção de uma vasta gama de metabólitos polares, apolares ou voláteis. Essas tecnologias são essenciais para monitorar intervenções terapêuticas, estudos de microbioma e alterações em vias metabólicas associadas a diferentes estados biológicos \cite{pang2024metaboanalyst}.

O processo de geração de dados segue um fluxo de trabalho rigoroso que começa com o design experimental, fundamental para minimizar vieses e variações aleatórias por meio da randomização. Após a coleta e preparação de amostras, a aquisição de dados brutos gera sinais complexos (espectros ou cromatogramas) que precisam de pré-processamento \cite{LILAND2011827,anwardeen2023statistical}. Esta etapa crítica envolve a redução de ruído, correção de linha de base, detecção de picos (peak picking) e alinhamento de picos entre as amostras para converter sinais brutos em matrizes de intensidade interpretáveis \cite{sun2024pretreating}.

Entretanto, a geração desses volumes massivos de dados traz desafios significativos. Os conjuntos de dados metabolômicos apresentam alta dimensionalidade, em que o número de variáveis frequentemente excede o número de amostras ($N \ll P$), o que complica a análise estatística tradicional \cite{LILAND2011827,picone2026integrating}. Além disso, os dados são marcados pela esparsidade e presença de valores ausentes, causados por metabólitos abaixo do limite de detecção ou falhas técnicas de alinhamento \cite{li2016performance}. Outro fator complicador são os efeitos de lote (batch effects) e a deriva instrumental, que podem introduzir variações não biológicas e obscurecer o sinal biológico real \cite{li2016performance}.

É importante destacar que, mesmo após o pré-processamento, podem permanecer falsos positivos na matriz final — valores que não representam um metabólito verdadeiro, mas sim ruídos provenientes do equipamento ou do ambiente \cite{schiffman2019filtering}. Para mitigar esse risco, softwares especializados como o \textbf{Progenesis QI} incorporam rotinas adicionais de controle de qualidade, como a redução de dados sem perda de informação essencial, o mascaramento de regiões cromatográficas ruidosas, alertas automáticos de inconsistências de calibração e métricas visuais de alinhamento. Além disso, o software oferece funções de normalização robusta e ferramentas de estatística multivariada (como PCA e dendrogramas) que auxiliam na detecção de outliers e na validação visual dos sinais \cite{watersprogenesisguide}.

Devido a essa complexidade, o tratamento adequado das informações tornou-se uma das etapas mais desafiadoras e demoradas. O avanço das tecnologias aumentou a necessidade de ferramentas computacionais, softwares especializados e abordagens in silico capazes de auxiliar no processamento, anotação e visualização de dados multivariados. Assim, a integração de métodos de aprendizado de máquina e estatística avançada é indispensável para extrair padrões significativos e garantir a reprodutibilidade dos resultados em experimentos de larga escala.


\section{\textbf{Cromatografia Líquida acoplada à Espectrometria de Massas (LC-MS)}}

A cromatografia líquida acoplada à espectrometria de massas (LC-MS) consolidou-se como uma das principais técnicas analíticas utilizadas em metabolômica, especialmente em abordagens não direcionadas. Essa metodologia combina duas etapas fundamentais: a separação cromatográfica dos compostos presentes na amostra e a detecção desses compostos pelo espectrômetro de massas. A cromatografia líquida permite separar metabólitos de acordo com suas propriedades físico-químicas, como polaridade, hidrofobicidade, carga e interação com a coluna cromatográfica, o que é essencial para lidar com a elevada complexidade química de matrizes biológicas como plasma, soro, urina e extratos celulares \cite{chen2022guide}.

Após a separação, os compostos são direcionados ao espectrômetro de massas, onde são ionizados e analisados de acordo com sua razão massa/carga (m/z). Essa etapa fornece informações cruciais, como massa acurada, padrão isotópico e fragmentação molecular, permitindo a caracterização de metabólitos em baixas concentrações. Entretanto, durante o processo de ionização, um mesmo metabólito pode originar diferentes sinais relacionados, como adutos, fragmentos ou formas isotópicas, aumentando a complexidade da interpretação dos dados \cite{mahieu2017systems}.

Em equipamentos de alta resolução (LC-HRMS), é possível obter maior precisão na determinação da massa dos compostos, o que contribui para a anotação putativa dos metabólitos. No entanto, essa identificação deve ser interpretada com cautela, já que compostos distintos podem apresentar massas semelhantes, padrões de fragmentação parecidos ou formar adutos e isômeros relacionados. Em estudos de metabolômica não direcionada, cada sinal detectado é geralmente definido por três parâmetros principais: tempo de retenção, valor de m/z e intensidade. Esses sinais, chamados de features metabolômicas, são organizados em uma matriz de dados que servirá de base para etapas subsequentes de pré-processamento, normalização e análise estatística \cite{schrimpe2016untargeted}.

A LC-MS, portanto, oferece elevada sensibilidade e ampla cobertura metabólica, sendo capaz de detectar simultaneamente centenas ou milhares de metabólitos em uma única análise. Essa capacidade amplia significativamente o potencial de descoberta biológica, mas também aumenta os desafios computacionais e estatísticos envolvidos no tratamento, integração e interpretação das informações metabolômicas. Dessa forma, compreender as particularidades da LC-MS é essencial para a escolha adequada de técnicas de pré-processamento e controle de qualidade, garantindo maior confiabilidade e robustez nos resultados obtidos.


\section{\textbf{Características dos Dados Metabolômicos}}

Os dados provenientes de experimentos metabolômicos apresentam características específicas que tornam sua análise significativamente mais complexa quando comparada a conjuntos de dados tradicionais. Essas características estão relacionadas tanto à natureza biológica dos metabólitos quanto às limitações e particularidades das técnicas analíticas utilizadas durante a aquisição dos dados. A metabolômica envolve a caracterização em larga escala de metabólitos presentes em células, tecidos e organismos, utilizando principalmente plataformas analíticas como espectrometria de massas (MS) e ressonância magnética nuclear (NMR), cada uma com vantagens, limitações e requisitos específicos de processamento e interpretação dos dados \cite{chen2022guide}.

Em geral, os datasets metabolômicos são caracterizados por alta dimensionalidade, elevada variabilidade experimental, presença de ruídos, valores ausentes e grande heterogeneidade entre amostras. Além disso, os dados frequentemente apresentam distribuições assimétricas, relações não lineares entre variáveis, diferenças expressivas de escala entre variáveis e forte sensibilidade às condições experimentais. A identificação e quantificação dos metabólitos constituem tarefas complexas, que não podem ser completamente automatizadas e exigem cuidadoso processamento computacional e análise estatística para obtenção de informações confiáveis \cite{shahjaman2021rmisbeta}.

Outro aspecto relevante é que pequenas alterações instrumentais ou diferenças no preparo das amostras podem gerar impactos significativos nos sinais medidos, dificultando a separação entre variações técnicas e variações biológicas reais. Fatores como deriva instrumental, diferenças entre lotes experimentais, sensibilidade analítica e limitações das plataformas de aquisição podem introduzir padrões artificiais nos dados, tornando necessárias múltiplas etapas de pré-processamento, normalização e controle de qualidade antes da realização das análises estatísticas \cite{shahjaman2021rmisbeta}.

A complexidade desses dados também está diretamente relacionada à capacidade das tecnologias modernas, como espectrometria de massas (MS) e ressonância magnética nuclear (NMR), de detectar simultaneamente centenas ou milhares de metabólitos em uma única análise. Enquanto técnicas baseadas em MS oferecem alta sensibilidade e ampla cobertura metabólica, métodos baseados em NMR fornecem elevada reprodutibilidade e informações estruturais detalhadas dos compostos analisados. Embora essas tecnologias ampliem significativamente o potencial de descoberta biológica, também aumentam os desafios computacionais e estatísticos envolvidos no tratamento, integração e interpretação das informações metabolômicas \cite{NMR}.

Dessa forma, compreender as características dos dados metabolômicos é essencial para a escolha adequada de técnicas de pré-processamento, normalização e análise estatística, garantindo maior confiabilidade, robustez e interpretabilidade dos resultados obtidos.

\subsection{Alta Dimensionalidade}

Os dados metabolômicos são caracterizados por alta dimensionalidade, situação em que o número de variáveis analisadas frequentemente excede o número de amostras disponíveis, configurando o problema conhecido como $p \gg n$ \cite{gunawan2022data} . Nesse cenário, cada metabólito detectado pelas plataformas analíticas corresponde a uma variável do conjunto de dados, fazendo com que experimentos metabolômicos possam gerar centenas ou milhares de variáveis simultaneamente, mesmo em estudos com número limitado de amostras biológicas.

Essa característica está diretamente relacionada ao avanço de tecnologias analíticas modernas, como espectrometria de massas (MS) e ressonância magnética nuclear (NMR), capazes de detectar grande quantidade de compostos metabólicos em uma única análise \cite{NMR}. Como consequência, os datasets metabolômicos tornam-se altamente complexos, multidimensionais e frequentemente acompanhados de ruídos experimentais, variabilidade técnica e diferenças de escala entre variáveis.

Em geral, os dados metabolômicos são organizados em matrizes bidimensionais, nas quais as linhas representam as amostras biológicas e as colunas correspondem às intensidades dos metabólitos detectados. Essa estrutura matricial constitui a base para aplicação de métodos estatísticos, técnicas de redução de dimensionalidade e algoritmos de aprendizado de máquina \cite{gunawan2022data}.

No entanto, a elevada dimensionalidade dos dados impõe desafios significativos à análise estatística e computacional \cite{shahjaman2021rmisbeta}. Métodos estatísticos tradicionais podem apresentar baixo desempenho nesse contexto, especialmente devido ao aumento do risco de overfitting, à presença de correlações espúrias e à redução da capacidade de generalização dos modelos analíticos. Além disso, o aumento do número de variáveis dificulta a identificação de padrões relevantes, fenômeno frequentemente associado à chamada maldição da dimensionalidade \cite{gunawan2022data}.

Diante desse cenário, torna-se necessária a aplicação de técnicas específicas de pré-processamento e redução de dimensionalidade, como transformação de dados, seleção de variáveis e métodos como a Análise de Componentes Principais (PCA). Essas abordagens buscam reduzir a complexidade dos datasets metabolômicos, preservar informações biologicamente relevantes e aumentar a robustez das análises estatísticas e computacionais realizadas.

\subsection{Variabilidade Técnica e Biológica}

A metabolômica, com seus dados de alta dimensionalidade, enfrenta o desafio de separar os efeitos mistos que influenciam os níveis de metabólitos. Muitos métodos estatísticos e experimentais estruturados, como design de experimentos e modelagem multivariada, permitem decompor a variabilidade e associar cada metabólito às causas específicas. Esses métodos ajudam a identificar assinaturas fisiológicas ou patológicas de forma mais confiável \cite{boccard2023we}.

Para a identificação dessas assinaturas, são utilizadas estratégias como o uso de razões entre metabólitos (normalização interna e melhor interpretação de vias metabólicas), modelagem de redes metabólicas (conecta substratos e produtos para entender as rotas bioquímicas) e estudos longitudinais (acompanham alterações ao longo do tempo). Essas estratégias são essenciais para transformar a metabolômica de uma simples busca por biomarcadores em uma abordagem mecanística, capaz de explicar como os processos bioquímicos mantêm seu estado de condições internas relativamente constantes, a homeostase \cite{boccard2023we}.

Em alguns estudos de metabolômica em larga escala, a coleta e análise de amostras podem durar dias, meses ou até anos. Esse período pode gerar os efeitos de lote (batch effects), que são variações artificiais nos dados causadas por diferenças no estado do instrumento, condições ambientais ou operadores. Os batch effects podem mascarar as verdadeiras relações biológicas entre metabólitos \cite{guo2023concordance}. 

São utilizados alguns métodos tradicionais para a correção de lote (BEC), que empregam amostras de controle de qualidade (QC), mas estas podem se deteriorar em experimentos longos. Outra alternativa são padrões internos isotopicamente marcados, que não cobrem bem todo o metaboloma. Com isso, métodos data-driven (baseados apenas nos dados) surgiram como solução \cite{guo2023concordance}.

Um método proposto no artigo \cite{guo2023concordance} é o CordBat, que seleciona um lote de referência como baseline, constrói um modelo gráfico gaussiano com dados do lote de referência e dos demais, e ajusta os coeficientes dos outros lotes para que suas correlações entre metabólitos fiquem em concordância com as do lote de referência.

As plataformas analíticas modernas, como LC-MS, GC-MS, NMR e ion mobility, geram dados altamente complexos e multivariados. Métodos estatísticos tradicionais muitas vezes não conseguem lidar com essa complexidade. O machine learning surgiu como uma solução, permitindo reconhecimento de padrões, classificação, regressão e seleção de variáveis em dados de alta dimensionalidade \cite{picone2026integrating}. 

Diante disso, a complexidade e variabilidade dos dados são fatores preocupantes para a utilização dos métodos tradicionais atualmente, visto a alta demanda de tempo necessário para análise e tratamento desses dados. Tecnologias como o aprendizado de máquina (machine learning), subcampo da Inteligência Artificial (IA), surgem com o intuito de otimizar e dar suporte aos processos realizados por equipes biomédicas.


\subsection{Valores ausentes e zero}

 Na metabolômica baseada em espectrometria de massas, o tratamento dos dados é um passo crítico para garantir qualidade e confiabilidade das análises. Durante a aquisição e pré-processamento, etapas como compressão de dados, filtragem de ruído e correção de baseline são fundamentais para melhorar a consistência dos sinais. Em seguida, procedimentos como detecção de picos, deconvolução e alinhamento permitem identificar corretamente os metabólitos presentes.

Um dos maiores desafios está relacionado aos valores ausentes e zeros. Esses podem ocorrer por diferentes razões: limites de detecção do equipamento, supressão iônica, desalinhamento de picos ou inconsistências na aquisição do sinal. A ausência de valores pode ser classificada como missing completely at random (MCAR), missing at random (MAR) ou missing not at random (MNAR), sendo este último frequentemente associado a metabólitos em concentrações abaixo do limite de detecção. Já a esparsidade descreve a presença de grande proporção de valores nulos ou próximos de zero, comum em análises não direcionadas, onde nem todos os metabólitos são detectados em todas as amostras \cite{metabolomicsdatatreatment}.

Esses problemas afetam diretamente a inferência estatística e o desempenho dos modelos de machine learning, podendo introduzir vieses e distorcer a estrutura de variância dos dados. Para mitigar tais efeitos, são aplicadas estratégias de pré-processamento como imputação de dados (k-nearest neighbors, imputação múltipla ou métodos baseados em ML, como random forest imputation), além de filtragem de variáveis e transformações matemáticas. O uso adequado dessas técnicas é essencial para preservar a relevância biológica, melhorar a robustez dos modelos e garantir interpretações confiáveis \cite{metabolomicsdatatreatment}.


\section{\textbf{Pré-processamento de Dados em Metabolômica}}

Na metabolômica, existe uma etapa denominada pré-processamento de dados, na qual diferentes procedimentos são aplicados aos dados brutos antes da realização de análises estatísticas e computacionais. Essa fase possui elevada complexidade, devido à necessidade de transformar grandes volumes de dados em informações consistentes e adequadas para interpretação analítica. Os dados vêm, em geral, de duas plataformas de aquisição, a espectroscopia por ressonância magnética nuclear (RMN) e a espectrometria de massas (EM), esta frequentemente associada a métodos de separação cromatográfica, e cada uma exige procedimentos próprios de pré-processamento \cite{LILAND2011827}. 

Em experimentos metabolômicos, normalmente é gerada uma grande quantidade de dados, os quais raramente podem ser utilizados de forma direta, completa e isenta de interferências experimentais. Esses dados frequentemente apresentam ruídos, variações instrumentais, desalinhamentos e outros fatores que podem comprometer significativamente a qualidade e a confiabilidade das análises subsequentes. Dessa forma, torna-se necessária a aplicação de técnicas de pré-processamento e pré-tratamento para adequar os dados às etapas analíticas posteriores. Na literatura metabolômica, os termos “pré-processamento de dados” e “pré-tratamento de dados” são frequentemente utilizados de maneira inconsistente \cite{sun2024pretreating}. De forma mais específica, o pré-processamento está relacionado às etapas iniciais aplicadas aos dados brutos, incluindo filtragem de ruído, correção de linha de base, faseamento, alinhamento de picos e binning espectral \cite{sun2024pretreating}. Em alguns trabalhos, procedimentos como normalização e dimensionamento de variáveis também são incorporados a essa etapa. Já o pré-tratamento é geralmente associado às técnicas de transformação, centralização e escalonamento dos dados, podendo incluir também métodos de estimativa ou imputação de valores ausentes.

De maneira geral, todas as etapas realizadas antes da análise estatística podem ser consideradas parte do processo de pré-processamento e pré-tratamento de dados. Entre as ciências ômicas, a metabolômica não direcionada destaca-se na busca por biomarcadores em sistemas biológicos complexos, permitindo analisar a resposta de organismos, tecidos ou células a estímulos externos e relacionar condições experimentais a alterações fenotípicas de interesse. Sua principal vantagem consiste na capacidade de fornecer um “instantâneo metabólico” do sistema biológico analisado. Entretanto, um dos principais desafios dessa abordagem está relacionado ao processamento de grandes volumes de dados, exigindo métodos computacionais robustos e eficientes para tratamento e interpretação das informações obtidas \cite{schrimpe2016untargeted}. Nesse contexto, as etapas de pré-processamento e pré-tratamento possuem papel fundamental, pois buscam minimizar artefatos instrumentais e variabilidades biológicas irrelevantes, melhorar a relação sinal-ruído (SNR, do inglês Signal-to-Noise Ratio) e transformar os dados em perfis espectrais mais interpretáveis, reduzindo sua dimensionalidade e aumentando a confiabilidade das análises metabolômicas \cite{sun2024pretreating}.

\subsection{Etapas de pré-processamento}

As etapas de pré-processamento no Progenesis QI seguem um fluxo de trabalho integrado projetado para preparar os dados de LC-MS para análise quantitativa e identificação posterior. O processo inicia-se com a Importação de Dados, fase na qual o software realiza uma rotina de redução de dados e aplica a calibração de massa (lock mass) para garantir a integridade e precisão das medições \cite{watersprogenesisguide}. Durante essa carga inicial, é realizada uma revisão de controle de qualidade (QC) para alertar sobre possíveis erros, como polaridade incorreta ou arquivos convertidos indevidamente para o formato centroided. Também é possível utilizar máscaras para excluir áreas ruidosas do mapa de íons, geralmente no início ou fim da corrida, causadas pela regeneração da coluna \cite{watersprogenesisguide}.

Em seguida, o Alinhamento atua compensando variações nos tempos de retenção entre as diferentes corridas analíticas, assegurando que as mesmas características sejam comparadas de forma consistente em todo o conjunto de dados. Esse processo exige a escolha de uma ``corrida de referência'', que pode ser selecionada automaticamente pelo software através da análise de similaridade entre todas as amostras, ou manualmente pelo usuário \cite{watersprogenesisguide}. O software fornece métricas visuais de qualidade, utilizando cores que variam do verde ao vermelho para indicar a precisão do alinhamento em diferentes regiões do mapa de intensidade de íons \cite{watersprogenesisguide}.

A detecção propriamente dita ocorre na etapa de Peak Picking, na qual são estabelecidos parâmetros e limiares para localizar os picos dos íons dos compostos em todas as corridas alinhadas \cite{watersprogenesisguide}. O usuário pode ajustar a sensibilidade da detecção e definir janelas mínimas para a largura de pico cromatográfico \cite{watersprogenesisguide}. Uma funcionalidade importante permite que apenas corridas de alta qualidade contribuam para o padrão de detecção de picos, evitando que amostras ruidosas prejudiquem o mapeamento de características em todo o experimento \cite{watersprogenesisguide}.

Posteriormente, a Normalização ajusta as abundâncias para permitir comparações estatísticas justas entre as amostras. O sistema identifica automaticamente a corrida ``menos diferente'' para servir como referência de normalização \cite{watersprogenesisguide}. Embora o método padrão considere todos os compostos detectados, o software oferece flexibilidade para normalizar com base em padrões externos ou em um conjunto específico de compostos de referência (housekeeping). Na determinação do fator de normalização, compostos que se comportam como outliers são automaticamente ignorados pelo algoritmo \cite{watersprogenesisguide}.

Por fim, a Desconvolução agrupa automaticamente as diversas formas de adutos de um único composto, consolidando a quantificação precisa da molécula original a partir da sua massa neutra \cite{watersprogenesisguide}. O cálculo da massa neutra depende da detecção de pelo menos dois adutos diferentes para o mesmo composto. Esta fase permite uma revisão detalhada, na qual o analista pode adicionar ou remover manualmente adutos específicos caso seus perfis cromatográficos ou de espectro de massa não correspondam à maioria dos íons agrupados para aquele composto \cite{watersprogenesisguide}.

\subsection{O papel do Progenesis QI no Pré-processamento Primário }

No âmbito do pré-processamento de dados metabolômicos, o software Progenesis QI (proprietário da empresa estadunidense Waters) consolida-se como uma plataforma de referência para a extração e o tratamento primário de características analíticas \cite{watersprogenesisguide}. Suas funções mais relevantes concentram-se no pipeline inicial de processamento: o alinhamento cromatográfico para a correção de desvios nos tempos de retenção, a co-detecção de picos — que mitiga a ocorrência de dados ausentes através da integração do ruído de fundo — e a deconvolução espectral, responsável por agrupar fragmentos, adutos e isótopos em entidades moleculares únicas. Após a finalização desse processo de deconvolução e quantificação dos compostos, antecedendo qualquer tentativa de modelagem ou filtragem estatística multivariada no ambiente nativo do software, as informações semiprocessadas são exportadas, e o pré-processamento é completado na ferramenta proposta. Nesse ponto, a matriz de abundâncias biológicas encontra-se unificada em um formato estruturado, porém altamente suscetível a falsos positivos latentes e derivas instrumentais, configurando o cenário e o insumo ideais para que a ferramenta assuma o fluxo e inicie o refino quimiométrico assistido. 

\section{\textbf{Tratamento de Dados Metabolômicos}}

O tratamento de dados metabolômicos constitui uma etapa crítica para transformar sinais analíticos brutos em matrizes limpas e comparáveis, aptas para análises estatísticas e interpretações biológicas \cite{datanor,sun2024pretreating}. Esse processo é multifacetado e busca minimizar artefatos instrumentais, ruídos e variações irrelevantes, sendo geralmente estruturado em fases de pré-processamento, correção de valores ausentes, normalização, escalonamento e transformação.

No pré-processamento inicial, tanto em espectrometria de massas (MS) quanto em ressonância magnética nuclear (NMR), são aplicadas técnicas para garantir a qualidade do sinal e a extração de características relevantes. Em MS, destacam-se a redução de ruído, correção de linha de base, detecção de picos e alinhamento cromatográfico para mitigar variações de tempo de retenção \cite{sun2024pretreating}. Já em NMR, são comuns procedimentos como correção de fase, remoção de solvente e alinhamento de picos, utilizando frequentemente algoritmos como o icoshift \cite{LILAND2011827}. Estratégias como o binning também são utilizadas para reduzir a dimensionalidade e suavizar ruídos, embora possam acarretar perda de informação em casos de desalinhamento severo.

A presença de valores ausentes e outliers é um desafio recorrente que pode comprometer as correlações entre variáveis. Esses valores faltantes podem decorrer de limites de detecção, supressão iônica ou inconsistências técnicas no alinhamento de picos \cite{sun2024pretreating}. Para lidar com isso, são aplicadas técnicas de imputação, como o k-nearest neighbors (k-NN), imputação múltipla ou métodos baseados em aprendizado de máquina, como random forest imputation. É prática comum descartar metabólitos com alta proporção de valores ausentes, seguindo a chamada “regra dos 80\%”, que retém variáveis presentes em pelo menos 80\% das amostras de algum grupo \cite{sun2024pretreating}. Outliers, por sua vez, podem ser identificados em análises exploratórias como PCA ou por medidas estatísticas como a distância de Hotelling’s T2 \cite{schiffman2018data}.

A normalização é considerada uma das etapas mais críticas do fluxo de trabalho, pois busca remover variações sistemáticas indesejadas e tornar as amostras comparáveis \cite{sun2024pretreating}. Métodos dirigidos por dados incluem a normalização pela soma total (TIC), mediana ou quantis. A Probabilistic Quotient Normalization (PQN) é amplamente recomendada em estudos de biofluidos por sua robustez ao utilizar o quociente mais provável em relação a um espectro de referência \cite{picone2026integrating}. Abordagens baseadas em padrões internos ou amostras de controle de qualidade (QC) são essenciais para corrigir a deriva instrumental ao longo de grandes lotes \cite{datanor}. Embora às vezes sejam agrupadas com a normalização, elas formam uma etapa própria, anterior a ela, a correção do sinal, e é assim que este trabalho as trata. Na correção por regressão local sobre os QCs (QC-RLSC), ajusta-se, para cada variável, uma curva LOESS da intensidade dos QCs em função da ordem de injeção, e cada amostra é corrigida pelo valor da curva no seu ponto de injeção \cite{dunn2011procedures}. Algoritmos mais recentes, como o Systematic Error Removal using Random Forest (SERRF), foram propostos para reduzir erros técnicos por meio da exploração de correlações entre compostos \cite{fan2019systematic}.

As etapas de escalonamento (scaling) e transformação ajustam a variância das variáveis, tornando metabólitos de diferentes ordens de abundância comparáveis entre si \cite{sun2024pretreating}. O autoscaling confere peso igual a todas as variáveis, mas pode inflar ruídos instrumentais. O Pareto scaling surge como uma alternativa intermediária popular, escalando os dados pela raiz quadrada do desvio padrão para reduzir a influência de metabólitos muito abundantes, mantendo os dados mais próximos da estrutura original do que o autoscaling \cite{mayonu2026evaluation}. Transformações logarítmicas e de potência são aplicadas para estabilizar variâncias e reduzir a assimetria característica de dados metabolômicos. Métodos como a Variance Stabilization Normalization (VSN) combinam a normalização entre amostras com uma estabilização não linear da variância, sendo ideais para análises exploratórias \cite{sun2024pretreating}.

Sendo assim, o tratamento de dados metabolômicos deve ser cuidadosamente planejado de acordo com a pergunta biológica, a plataforma analítica utilizada (MS ou NMR) e a sensibilidade dos métodos estatísticos subsequentes. A escolha adequada de cada etapa garante maior robustez, precisão e interpretabilidade dos resultados, transformando dados brutos complexos em informações confiáveis para estudos clínicos, ambientais e de segurança alimentar \cite{picone2026integrating}.

\section{\textbf{Normalização de Dados}}

A normalização dos dados é uma etapa crucial e extremamente crítica nos estudos de metabolômica, especialmente quando os dados vêm de plataformas de espectrometria de massas. O objetivo é reduzir a variabilidade quantitativa e garantir análises estatísticas robustas e confiáveis \cite{low2026multifaceted}. Sem essa etapa, os resultados podem refletir ruídos técnicos ou ambientais em vez de variações biológicas reais, aumentando o risco de falsos positivos e comprometendo a interpretação dos achados.

A normalização em metabolômica é multifacetada e envolve três dimensões principais: normalização de amostras, correção de sinais e transformações estatísticas. Cada uma atende a necessidades específicas, mas juntas garantem maior precisão analítica e interpretabilidade dos resultados. O artigo \cite{low2026multifaceted} enfatiza que a escolha do método deve considerar as particularidades da metabolômica, como a presença de valores ausentes e a alta variabilidade técnica, e recomenda estratégias rigorosas para avaliar os resultados da normalização. Assim, a normalização é essencial para reduzir ruídos, aumentar a confiabilidade das análises e permitir que os dados reflitam de forma mais fiel os processos biológicos.

Contudo, a diversidade de métodos disponíveis torna a normalização uma tarefa complexa e muitas vezes considerada uma “ciência imprecisa”. Para enfrentar esse desafio, foi desenvolvido o pacote NormalizeMets, em R com interface em Excel, que permite a avaliação comparativa de diferentes métodos de normalização. Essa ferramenta auxilia pesquisadores a avaliar, selecionar e implementar o método mais adequado ao seu conjunto de dados e objetivo de pesquisa. Além disso, o NormalizeMets oferece recursos de visualização interativa e análises estatísticas finais, como clustering, classificação, identificação de biomarcadores e correlação, tornando-se uma solução prática e acessível para diferentes perfis de usuários \cite{de2018normalizemets}.

Portanto, a normalização em metabolômica não deve ser vista apenas como uma etapa técnica, mas como um processo estratégico que impacta diretamente a qualidade dos resultados e a validade das interpretações biológicas.

\subsection{Tipos de Normalização}

\subsubsection{Métodos baseados em dados}

Os métodos baseados em dados utilizam apenas a matriz de valores obtida nos experimentos para reduzir variabilidade e corrigir distorções. São estratégias como log-transformação, normalização pela soma, média ou mediana, TIC (total ion chromatogram), PQN (probabilistic quotient normalization), VSN (variance stabilizing normalization) e LOESS. Esses métodos são amplamente aplicados em estudos de metabolômica, especialmente em análises não direcionadas, por não dependerem de padrões externos. Contudo, diferentes métodos podem influenciar diretamente a descoberta de biomarcadores e a interpretação dos resultados, sendo necessário avaliar cuidadosamente qual técnica aplicar \cite{datanor,li2016performance}.

\subsubsection{Métodos baseados em padrões internos}

Os métodos baseados em padrões internos utilizam compostos adicionados às amostras como referência para corrigir variações técnicas. Podem ser padrões marcados com isótopos estáveis (SIL-IS) ou compostos estruturalmente semelhantes. Essa abordagem é comum em análises de LC-MS e LC-HR-MS/MS, especialmente em lipidômica, pois permite ajustar variações de resposta do instrumento e melhorar a precisão quantitativa. A escolha adequada dos padrões internos é crítica, já que nem sempre há correlação linear entre o padrão e o metabólito alvo. Diretrizes recentes recomendam o uso de padrões representativos de diferentes classes químicas para garantir maior cobertura e confiabilidade \cite{drotleff2019guidelines}, \cite{datanor}.

\subsubsection{Métodos baseados em QC}

Os métodos baseados em QC (Quality Control) utilizam amostras de controle de qualidade, geralmente preparadas a partir de alíquotas de todas as amostras experimentais, para monitorar estabilidade do instrumento e corrigir variações sistemáticas. Estratégias como pooled QC samples, referência material (Ref-M) e algoritmos avançados como SERRF (Systematic Error Removal using Random Forest) são exemplos. Esses métodos permitem reduzir erros técnicos, corrigir efeitos de lote e aumentar a comparabilidade entre diferentes conjuntos de dados. Além disso, o uso de QC é considerado essencial para validar a integridade dos dados e garantir que os resultados reflitam variações biológicas reais \cite{yao2022normalization}, \cite{fan2019systematic}, \cite{drotleff2019guidelines}.

\section{\textbf{Controle de Qualidade de Dados Metabolômicos}}

O controle de qualidade (QC) é uma etapa indispensável nos estudos de metabolômica, pois garante que os dados obtidos reflitam variações biológicas reais e não artefatos técnicos ou ambientais. A ausência de procedimentos robustos de QC pode resultar em interpretações equivocadas, aumento de falsos positivos e comprometimento da reprodutibilidade dos resultados \cite{datanor}. 

As amostras de QC são geralmente preparadas a partir de alíquotas de todas as amostras experimentais, formando um pooled sample que representa o perfil metabólico médio do estudo. Essas amostras são injetadas periodicamente durante a corrida analítica para monitorar a estabilidade do instrumento, identificar deriva de sinal e corrigir efeitos de lote (batch effects) \cite{fan2019systematic}. Além disso, materiais de referência (Ref-M) e padrões internos isotopicamente marcados podem ser utilizados para avaliar a consistência das medições e ajustar variações técnicas \cite{drotleff2019guidelines}.

O QC também desempenha papel fundamental na filtragem de features pouco reprodutíveis. Uma métrica amplamente utilizada é o desvio padrão relativo (RSD), que avalia a variabilidade de cada metabólito nas amostras de controle. Features com RSD elevado são consideradas instáveis e podem ser removidas da matriz de dados, aumentando a confiabilidade das análises subsequentes \cite{yao2022normalization}. Embora não exista um ponto de corte universal, valores de RSD inferiores a 20–30\% são frequentemente adotados como referência em estudos metabolômicos.

Recentemente, algoritmos avançados têm sido incorporados ao QC para aumentar sua eficiência. O Systematic Error Removal using Random Forest (SERRF) é um exemplo de abordagem baseada em aprendizado de máquina que utiliza correlações entre metabólitos para eliminar erros sistemáticos e corrigir deriva instrumental \cite{fan2019systematic}. Essa integração entre estatística tradicional e inteligência artificial fortalece o papel do QC como etapa estratégica, capaz de reduzir ruídos e aumentar a comparabilidade entre diferentes conjuntos de dados.

Portanto, o controle de qualidade em metabolômica não deve ser visto apenas como uma etapa auxiliar, mas como um componente central do fluxo analítico. Ele assegura que os resultados obtidos sejam estatisticamente consistentes, analiticamente confiáveis e biologicamente interpretáveis, criando a base necessária para análises robustas e reprodutíveis.


\section{\textbf{Análise Estatística de Dados Metabolômicos}}

A análise estatística é uma etapa central nos estudos metabolômicos, sendo responsável por transformar grandes volumes de dados brutos em informações interpretáveis. Devido à alta dimensionalidade, à forte correlação entre as variáveis e à presença de ruídos, a escolha cuidadosa do método estatístico adequado é fundamental para garantir resultados confiáveis e evitar a perda de informações relevantes ou a identificação de tendências incorretas \cite{anwardeen2023statistical}.

De maneira geral, as abordagens estatísticas aplicadas em metabolômica podem ser divididas em dois grandes grupos: métodos univariados e métodos multivariados. Os métodos univariados, como o teste t de Student, ANOVA, teste de Mann-Whitney e correções para múltiplas comparações (Bonferroni, Benjamini-Hochberg), analisam cada metabólito de forma isolada, sendo amplamente utilizados pela facilidade de aplicação e interpretação. Já os métodos multivariados consideram simultaneamente o comportamento de todas as variáveis, permitindo identificar padrões, correlações e estruturas latentes nos dados \cite{anwardeen2023statistical,antonelli2019statistical}.

Dentre as técnicas multivariadas mais utilizadas em metabolômica, destacam-se a PCA (Principal Component Analysis), aplicada principalmente como ferramenta de exploração preliminar dos dados, e a PLS-DA (Partial Least Squares Discriminant Analysis), um dos métodos mais usados em estudos baseados em espectrometria de massas. Métodos baseados em ANOVA multivariada, como ASCA (ANOVA Simultaneous Component Analysis), rMANOVA (regularized MANOVA) e GASCA (Group-wise ANOVA Simultaneous Component Analysis), têm sido propostos como alternativas para decompor a variabilidade dos dados e associar metabólitos a fatores experimentais específicos, sendo particularmente úteis em estudos com delineamentos experimentais complexos baseados em LC-MS \cite{perezcova2022comparison}.

Mais recentemente, métodos de aprendizado de máquina (machine learning), como Random Forest (RF), Support Vector Machines (SVM), redes neurais artificiais e modelos esparsos como LASSO (Least Absolute Shrinkage and Selection Operator) e SPLS (Sparse Partial Least Squares), têm sido incorporados às pipelines de análise metabolômica \cite{anwardeen2023statistical}. Estudos comparativos demonstraram que, em cenários nos quais o número de metabólitos é semelhante ou superior ao número de amostras, comuns em análises não direcionadas com coortes pequenas, modelos multivariados esparsos apresentam maior robustez estatística, maior seletividade e menor potencial de identificação de relações espúrias quando comparados às abordagens univariadas tradicionais \cite{henglin2022quantitative}.

Além disso, ferramentas computacionais como o MetaboAnalyst têm se consolidado como plataformas amplamente utilizadas para análise estatística metabolômica, oferecendo módulos integrados para análises univariadas e multivariadas, classificação por RF, análise de biomarcadores, análise de enriquecimento e mapeamento de vias metabólicas. Essa integração possibilita uma análise abrangente desde o processamento inicial dos espectros até a interpretação funcional dos resultados \cite{pang2024metaboanalyst}.

\subsection{Métodos Exploratórios (PCA)}

A Análise de Componentes Principais (PCA) é amplamente reconhecida como um dos métodos exploratórios mais importantes na metabolômica, sendo classificada como uma técnica de aprendizado de máquina não supervisionado. Sua função principal é reduzir a dimensionalidade dos dados, transformando conjuntos complexos e de alta dimensão em componentes principais (PCs) que retêm a maior parte da variância original. Essa redução permite que os dados sejam representados em gráficos de dispersão bidimensionais ou tridimensionais, facilitando a visualização de padrões e agrupamentos entre amostras. Nos gráficos de PCA, os scores representam a distribuição das amostras e sua variância, enquanto os loadings indicam a contribuição de cada metabólito para a separação observada, auxiliando na interpretação biológica dos resultados \cite{mayonu2026evaluation}.

Uma das principais vantagens do PCA é o baixo risco de sobreajuste (overfitting), já que não depende de rótulos pré-definidos, diferentemente de métodos supervisionados como o PLS-DA \cite{jiang2022principal}. Isso o torna uma ferramenta robusta para análises exploratórias iniciais. Por outro lado, o PCA apresenta limitações quando há separação parcial entre grupos ou excesso de ruído nos dados, o que dificulta a interpretação visual e a identificação de metabólitos relevantes \cite{mayonu2026evaluation}. Além disso, em conjuntos de dados de alta dimensão, em que o número de variáveis ($p$) é muito maior que o de amostras ($n$), medidas tradicionais de detecção de outliers baseadas em PCA podem se tornar pouco confiáveis e aumentar a taxa de falsos positivos \cite{brini2021improved}.

Para superar essas limitações, novas abordagens têm sido propostas. Uma delas aplica valores projetados dos scores na direção da separação entre grupos, seguida de teste t para verificar estatisticamente a significância da separação. Outra estratégia combina os valores de loadings com testes estatísticos convencionais, como o t de Student ou Welch, aumentando a eficiência na descoberta de biomarcadores e metabólitos críticos \cite{mayonu2026evaluation}. O PCA também tem sido utilizado como etapa preliminar em modelos de classe única (one-class modeling), reduzindo dimensões antes da aplicação de medidas de distância, como a Distância de Mahalanobis, para detectar perfis metabólicos anormais em dados de alta dimensão \cite{brini2021improved}.

Em síntese, o PCA deve ser entendido não apenas como uma ferramenta exploratória visual, mas como um método que, quando combinado com testes estatísticos rigorosos e técnicas de projeção, torna-se robusto para a identificação de metabólitos importantes mesmo em condições de sobreposição de dados ou em estudos ambientais e clínicos complexos.


\section{\textbf{Desafios na Escolha de Métodos}}

A escolha do método de normalização em metabolômica é descrita como uma “área cinzenta” \cite{de2018normalizemets}, e a dos fluxos de trabalho analíticos e estatísticos em geral também é difícil, resultado da complexidade intrínseca dos dados e da ausência de protocolos padronizados \cite{wiedmer2013selection}. Um dos principais desafios é a alta dimensionalidade dos conjuntos de dados, marcada pelo problema “large $P$, small $n$” \cite{mayonu2026evaluation}, no qual o número de variáveis supera em muito o número de amostras. Essa condição compromete a aplicação de métodos tradicionais, como a MANOVA clássica, que exige mais amostras do que variáveis para ser estatisticamente viável, e aumenta significativamente o risco de sobreajuste em modelos supervisionados, como o PLS-DA \cite{perezcova2022comparison}.

A normalização, embora crítica, é frequentemente considerada imprecisa, uma vez que diferentes métodos podem alterar diretamente quais metabólitos são identificados como biomarcadores \cite{de2018normalizemets}. Além disso, ferramentas comerciais e plataformas de software oferecem algoritmos e coberturas distintas, o que dificulta a comparação e a harmonização de resultados entre diferentes estudos \cite{datanor,li2016performance}.

Outro ponto recorrente é a dificuldade em distinguir a variabilidade biológica da técnica. Fatores como dieta, estilo de vida, ritmo circadiano e manuseio de amostras introduzem ruídos e variações temporais que podem mascarar padrões biológicos sutis \cite{datanor,picone2026integrating}. Isso exige a aplicação de métodos robustos para correção de efeitos de lote (batch effects) e deriva instrumental de sinal \cite{fan2019systematic}. Pesquisadores destacam a importância do uso de materiais de referência e controles de qualidade (QCs) para monitorar e mitigar essas inconsistências \cite{datanor}.

No que diz respeito às plataformas analíticas, nenhum método isolado é capaz de cobrir todo o metaboloma devido à enorme diversidade química das moléculas \cite{wiedmer2013selection}. A espectrometria de massas (MS) oferece alta sensibilidade e seletividade, mas sofre com efeitos de matriz, supressão de íons e menor reprodutibilidade técnica \cite{elmsjo2018method}. Por outro lado, a ressonância magnética nuclear (NMR) é reconhecida por sua robustez e alta reprodutibilidade, porém apresenta sensibilidade limitada, o que pode ocultar metabólitos em baixas concentrações \cite{mayonu2026evaluation}. A complementaridade entre ambas as técnicas é, portanto, indispensável para uma visão holística do sistema biológico \cite{picone2026integrating}.

A ausência de protocolos padronizados permanece como uma barreira significativa. Muitas ferramentas computacionais possuem múltiplas configurações e parâmetros que raramente são reportados com clareza em publicações, permitindo que diferentes pesquisadores extraiam conclusões distintas de um mesmo conjunto de dados, introduzindo o chamado viés do usuário \cite{misra2016updates}.

Por fim, o desafio das “unknowns” (compostos desconhecidos) é central, especialmente em metabolômica não-alvo (untargeted). Grande parte dos sinais detectados por espectrometria de alta resolução não pode ser anotada com confiança devido a lacunas em bancos de dados espectrais e à escassez de padrões químicos autênticos para comparação, limitando a interpretação biológica e a descoberta de novos mecanismos de ação \cite{picone2026integrating}.

Em síntese, a escolha de métodos em metabolômica exige um equilíbrio rigoroso entre sensibilidade, robustez e interpretabilidade, sendo sempre condicionada pela pergunta biológica formulada e pelo desenho experimental adotado \cite{wiedmer2013selection}. Esses desafios justificam a necessidade de sistemas de apoio à decisão baseados em inteligência artificial, capazes de organizar fluxos de trabalho complexos, reduzir vieses e aumentar a reprodutibilidade das análises.

\section{\textbf{Avaliação da Qualidade dos Dados}}

A implementação de estratégias de controle de qualidade (QC) é crucial para assegurar a reprodutibilidade, a precisão e o significado biológico dos dados metabolômicos, minimizando erros humanos, aleatórios e sistemáticos \cite{QUINTAS2018137, doi:10.1021/acs.analchem.3c03660}. Uma avaliação robusta deve ser realizada de forma sequencial, abrangendo desde a exploração inicial dos dados brutos até a validação estatística final \cite{QUINTAS2018137, doi:10.1021/acs.analchem.3c03660}.

\subsection{Exploração Inicial e Estabilidade Instrumental}

O processo deve iniciar-se com a inspeção de amostras de branco (procedurais) para identificar sinais de artefatos, contaminantes ou arraste (carryover). Recomenda-se sinalizar picos como contaminantes quando seus valores médios nas amostras reais não excederem três vezes a média detectada nos brancos \cite{doi:10.1021/acs.analchem.3c03660}. Simultaneamente, a estabilidade do sistema é monitorada por meio de gráficos de série temporal e análises de componentes principais (PCA) para detectar derivas (drifts) de sinal relacionadas à ordem de injeção. Para a detecção de medições ``fora de controle'', utilizam-se as cartas de controle de Shewhart e as regras de Westgard, que permitem identificar falhas súbitas, como o entupimento de colunas, através da análise de desvios padrão (±SD) em compostos descritores químicos \cite{doi:10.1021/acs.analchem.3c03660}.

\subsection{Tratamento de Valores Faltantes e Outliers}

Uma abordagem moderna de QC diferencia valores faltantes (decorrentes de limitações técnicas) de dados verdadeiramente ausentes (devido a razões biológicas, como xenobióticos específicos de um grupo). Propõe-se que características moleculares com mais de 70\% de valores faltantes sejam tratadas como valores zero reais, enquanto para as demais variáveis utiliza-se a regra de filtragem de 20\% e posterior imputação \cite{doi:10.1021/acs.analchem.3c03660}. A detecção de outliers é realizada através do teste Hotelling T2 (elipse de 95\% no PCA) e do gráfico DmodX, que identifica observações com alta variância residual não explicada pelo modelo \cite{doi:10.1021/acs.analchem.3c03660}.

\subsection{Marcadores de Qualidade Pré-analítica}

Para mitigar o impacto de erros na coleta e processamento das amostras (fase pré-analítica), monitora-se um painel de metabólitos indicadores de integridade metabólica. Níveis anormais de hipoxantina, lactato e glicose podem indicar falhas no resfriamento (quenching) ou processamento tardio de amostras de sangue, enquanto a creatina e o ácido succínico são marcadores de crescimento bacteriano em urina \cite{doi:10.1021/acs.analchem.3c03660}.

\subsection{Comparabilidade e Correção de Efeito de Grupo}

Em estudos que envolvem a integração de múltiplos conjuntos de dados ou coortes longitudinais, a avaliação da qualidade deve considerar o ``efeito de grupo'' \cite{SALANON2025344753}. A estratégia PARSEC (Post-Acquisition corRection Strategy) permite harmonizar os dados sem depender de QCs de longo prazo, utilizando uma padronização por lote combinada com modelos de regressão mista. A eficiência desta correção é avaliada pela redução da variabilidade inter-grupo nos gráficos de PCA e pela queda na correlação entre os parâmetros de coleta (lote, coorte, ordem de injeção) e os componentes principais dos dados \cite{SALANON2025344753}.

\subsection{Métricas de Precisão e Exatidão}

A etapa final da avaliação consiste na aplicação de critérios de aceitação rigorosos para os descritores químicos:
Desvio Padrão Relativo (RSD) < 30\% para a intensidade do pico.
RSD < 2\% para o tempo de retenção (RT).
Erro de massa < 10 ppm (para espectrometria de alta resolução).
RSD < 15\% para a largura do pico.
Esses indicadores ajudam a verificar se a variação técnica não excede a variabilidade biológica interindividual, permitindo uma modelagem estatística confiável \cite{doi:10.1021/acs.analchem.3c03660, SALANON2025344753}. 

\section{\textbf{Limitações das Abordagens Atuais}}

Apesar dos avanços tecnológicos, a análise metabolômica enfrenta desafios significativos que limitam sua eficácia e reprodutibilidade global. As abordagens atuais lidam com barreiras que vão desde a variabilidade instrumental até a complexidade intrínseca dos sistemas biológicos.

\subsection{Reprodutibilidade e Variabilidade Instrumental}

Um dos principais obstáculos é a dificuldade em obter medições reprodutíveis ao longo do tempo e entre diferentes instrumentos. A acumulação de variâncias decorrentes da deriva instrumental (drift), instabilidade cromatográfica e mudanças nas condições ambientais pode alterar os resultados de forma substancial \cite{jong2018challenges}. Atualmente, não é trivial calibrar dados gerados em dias diferentes ou em plataformas distintas para alcançar medições comparáveis. Além disso, a supressão de íons e a dificuldade em replicar tempos de retenção entre sistemas diversos continuam sendo limitações críticas para a análise por LC-MS \cite{johnson2012challenges}.

\subsection{Desafios nos Padrões de Controle de Qualidade (QC)}

Embora padrões de QC sejam essenciais para medir efeitos de matriz, as soluções atuais apresentam limitações: 
\begin{itemize}
\item Amostras em pool: o uso de alíquotas agrupadas de todas as amostras do estudo é uma prática comum, mas torna-se impraticável para projetos de longo prazo devido à disponibilidade limitada de material \cite{jong2018challenges}.
\item Estabilidade e artefatos: amostras de referência (como plasma ou soro) exigem monitoramento rigoroso, pois as condições de armazenamento e a preparação podem gerar artefatos e variabilidade que confundem a identificação de metabólitos \cite{jong2018challenges, johnson2012challenges}.
\item Complexidade da matriz: a densidade de picos em matrizes biológicas torna irrealista o monitoramento de mais do que uma seleção limitada de compostos, apresentando muitos sinais interferentes para um controle de qualidade (QA/QC) verdadeiramente confiável \cite{jong2018challenges}.
\end{itemize}
\subsection{Identificação de Compostos e Lacunas nas Bases de Dados}

A identificação precisa de biomarcadores permanece como o maior gargalo da metabolômica. As principais dificuldades incluem:
\begin{itemize}
    \item Isômeros, isóbaros e fragmentos: múltiplos compostos biológicos podem compartilhar a mesma massa exata (isômeros, de mesma fórmula molecular) ou a mesma massa nominal (isóbaros), e compostos artefatuais ou fragmentos estruturalmente diferentes frequentemente superam em número os metabólitos conhecidos no estudo \cite{jong2018challenges}.
\end{itemize}

\begin{itemize}
    \item Bases de Dados Incompletas: Enquanto a GC-MS possui bibliotecas robustas, as bases de dados para LC-MS estão muito menos desenvolvidas, dificultando a comparação de dados entre laboratórios devido a diferenças em colunas, gradientes e energias de colisão \cite{johnson2012challenges}.
\end{itemize}

\subsection{Falta de Padronização e Compartilhamento de Dados }

Em 2012, apontava-se uma variação considerável nos métodos de amostragem, preparação e mineração de dados entre laboratórios e até entre pesquisadores de uma mesma instituição. A ausência, à época, de um banco de dados central de relatórios e de procedimentos operacionais padrão (SOPs) universais impedia a integração em larga escala de estudos metabolômicos e a validação cruzada de descobertas entre diferentes grupos de pesquisa \cite{johnson2012challenges}.